\chapter{Zusammenfassung und Auswertung}

\section{Vergleichende Auswertung der Clustering-Verfahren}

\section{Weitere Ansätze}

Der Vergleich der Algorithmen in dieser Arbeit entstand im Rahmen eines 
Praktikums und umfasst die zum damaligen Zeitpunkt bekannten und verfügbaren 
Verfahren. Im Nachgang wurde mit Brain \cite{brain} ein weiterer 
Log-Parsing-Algorithmus bekannt, der für den beschriebenen Anwendungsfall 
relevant ist. Da Brain nicht mehr in die systematische Evaluierungsphase 
eingeflossen ist, wird er an dieser Stelle ergänzend vorgestellt. Erste Tests 
mit denselben Datensätzen, die auch für den Vergleich der anderen Algorithmen 
verwendet wurden, deuten darauf hin, dass Brain sowohl Drain als auch den 
vektorbasierten Ansatz mit PostgreSQL und pg\_vector in der Parsing-Qualität 
übertrifft – und das bei deutlich geringerem Ressourcenverbrauch. Während der 
vektorbasierte Ansatz auf ein vortrainiertes Embedding-Modell sowie eine 
laufende Datenbankinstanz angewiesen ist, kommt Brain ohne diese Abhängigkeiten 
aus und arbeitet rein algorithmisch.
\section{Brain: Log Parsing mit bidirektionalem Parallelbaum}


\section{Nächste Schritte}

\section{Ausblick}
